\section{External measurements and a trajectory comparator}
We analyzed two published recordings using a protocol committed before scoring them. The first is the synchronous bead~1 force readout in two transverse directions from the passive C-Trap calibration released by Bugiel \cite{ctrap}. It contains 781,250 frames over ten seconds. The second is the 50,000-frame experimental position trace in the static harmonic thermophoretic feedback trap of Argun and colleagues \cite{argun}. The first fifth of each recording fixed the coordinate means and scales. In standardized coordinates the hub is the interval or square \([-0.75,0.75]^p\) in dimension \(p\). Each port starts beyond one face and, for \(p=2\), has transverse half-width 0.375. The corners fall in an exterior label. This calibration uses no held-out transitions. All subsequent measurements use the remaining 625,000 and 40,000 frames.

For each lag we divided directional pair counts by the number of available pairs, averaged each joint matrix with its transpose and used the resulting one-step row sum for \(M\). Symmetrization asks whether a reversible surrogate can reproduce these symmetric summaries. The archive preserves the directional counts and discrepancies. We repeated the comparison within each of eight contiguous blocks and after one shuffle preserving all label frequencies. The estimated flows and scores in overlapping windows are dependent. We report their plug-in state counts as descriptive sensitivity to lag and sampling.

\paragraph{A dependence-valid comparison with more information.}
A complete path records the port through which each visit enters the hub and the location of its next exit. For an entrance directly from port \(i\), record success if the first outside frame is again in \(i\). All other exits, including exits to the unmonitored exterior, are failures. Given any fine state \(h\) in the hub, let \(q_{hi}\) be its probability of first exiting into \(i\). At an entrance the conditional success probability given the observed past is at most \(b_i=\max_h q_{hi}\), by the strong Markov property. If \(b_i>1/2\) for \(r\) ports, the model needs at least \(r\) distinct hub states since \(\sum_i q_{hi}\le 1\). This argument also applies to nonreversible homogeneous models.

With \(S\) same-port exits out of \(n\) visits entered from port \(i\), we test \(b_i\le1/2\) using
\[
 L_i(n)=\tfrac14\sum_{q\in\{.55,.65,.75,.90\}}(2q)^S[2(1-q)]^{n-S}.
\]
At an eligible entrance each factor for a fixed \(q\) has conditional expectation at most one under this null. Skipping entrances from other ports leaves the product unchanged. Thus the mixture is a nonnegative supermartingale on the sequence of visits, starting at one. Ville's inequality bounds the chance of ever reaching 120 by \(1/120\) \cite{howard}. The final completed visits form a prefix even when the recording ends within the hub. A union bound over six ports in two records gives simultaneous 95\% lower bounds on their respective model dimensions. This inference assumes the held-out labels follow a homogeneous finite-state Markov representation at the acquisition frame rate and that the labels are deterministic functions of its states. It imposes no independence or mixing-rate assumption on successive frames. Electronic filtering, sensor memory and feedback can contribute to the state requirement of the observed process.

As a separate one-state check we compared the arrival port of one completed excursion outside \(B\) with the departure port of the next, using disjoint successive pairs. Under a one-state hub, excursions regenerate. Discard the initial and terminal partial excursions, self holds and excursions whose first or last outside label is unmonitored. The order of the remaining complete excursions is exchangeable under the null: permuting complete excursions preserves their product likelihood, total duration and both observation boundaries. Permuting the even excursions conditional on the two margins gives a right-tail test for excess port matches, with 9,999 random permutations and a plus-one correction. We apply Holm's adjustment across the two chronological records for this separate family.

\begin{table}[htbp]
\centering\small
\caption{Held-out endpoint scores and local state-count comparisons. For each recording the floor shown is the sharp floor for one fewer state than the number of ports. Count denotes the pooled plug-in lower bound; brackets give its range across eight contiguous blocks. Masses are normalized to total mass one.}
\label{tab:measured}
\begin{tabular}{r rrr rrr}
\toprule
 & \multicolumn{3}{c}{Passive calibration (four ports)} & \multicolumn{3}{c}{Harmonic trace (two ports)}\\
\cmidrule(lr){2-4}\cmidrule(lr){5-7}
Lag (frames) & \(Y\) & \(\E_3\) & Count & \(Y\) & \(\E_1\) & Count\\
\midrule
1    & .0000792 & .000110 & 4 [3,4] & 0       & .000223 & 2 [2,2]\\
4    & .00651   & .000684 & 1 [1,1] & 0       & .00120  & 2 [2,2]\\
16   & .0219    & .00133  & 1 [1,1] & .000363 & .00478  & 2 [2,2]\\
64   & .0259    & .00138  & 1 [1,1] & .0158   & .0167   & 2 [1,2]\\
256  & .0260    & .00135  & 1 [1,1] & .0405   & .0255   & 1 [1,1]\\
1024 & .0252    & .00133  & 1 [1,1] & .0394   & .0256   & 1 [1,2]\\
\bottomrule
\end{tabular}
\end{table}

In the passive calibration the four-state plug-in result occurred at one frame (12.8 microseconds) and in seven of its eight blocks. The pooled three-state margin was \(3.05\cdot10^{-5}\) in joint mass. For scale, the maximum forward--reverse discrepancy among two-step entries was \(1.44\cdot10^{-4}\). These statistics involve different entries; the discrepancy illustrates the resolution required to interpret the margin. At four frames the scalar already permitted a one-state hub. Both shuffled recordings returned only the trivial endpoint count at every lag. The exact and occupancy-relaxed bounds gave identical pooled counts at all twelve measured settings. In particular the relaxed three-state floor at one frame was \(1.05\cdot10^{-4}\), also above the observed mass. The benefit of sharpness in these data is a larger numerical margin; its decision advantage appears in the constructive example and one grid configuration.

The separate return-port test found 1,705 matches among 4,914 excursion pairs in the passive recording (conditional null expectation 1,232.0), and 385 among 453 in the harmonic recording (expectation 226.6). Each raw permutation \(p\)-value was 0.0001, giving 0.0002 after Holm adjustment. The corresponding shuffled \(p\)-values were 0.546 and 0.872. The port-specific supermartingale test supported at least two states in the harmonic recording (381/445 and 399/462 same-port exits). In the passive recording all four fractions lay between 0.268 and 0.294, yielding the trivial lower bound of one by that test. Thus full-path evidence excludes a one-state homogeneous hub in both experiments; the four-state figure for the passive record remains a descriptive two-lag comparison. All results concern memory in a model of the digitized coordinates, which include the instrument response.

